\documentclass[11pt,a4paper]{article}
\usepackage[utf8]{inputenc}
\usepackage[T1]{fontenc}
\usepackage{lmodern,amsmath,amssymb,amsthm,mathtools,booktabs,hyperref,enumitem,microtype}
\usepackage[margin=1in]{geometry}
\usepackage[numbers,sort&compress]{natbib}

\newtheorem{definition}{Definition}[section]
\newtheorem{proposition}[definition]{Proposition}
\newtheorem{hypothesis}[definition]{Hypothesis}
\newtheorem{corollary}[definition]{Corollary}
\theoremstyle{remark}
\newtheorem*{remark}{Remark}

\newcommand{\Frame}{\mathcal F}
\newcommand{\Seldon}{\mathcal S}
\newcommand{\Narr}{\mathcal N}
\newcommand{\Hist}{\mathcal H}

\title{The Law of Narrative Momentum --- Theory\\\large First-Class History, Directional Persistence, and Revision}
\author{Matthew Parslow\\Independent Researcher}
\date{}

\begin{document}
\maketitle

\begin{abstract}
This paper develops Narrative as a real, measurable, frame-relative quantity rather than as a metaphorical story stored inside the graph it describes. Graphic Equalisation establishes that a containing frame may possess first-class relational state about a lower graph even where that quantity is absent from the lower graph itself. Realised history can therefore exist as ordinary state in the representing frame.

Narrative is defined as a discrete frame-level measure of realised-history structure that remains operative in the present representation of a graph. The Narrative dimension need not exist in every frame; where it is absent, Narrative is undefined rather than numerically zero. Where history-bearing structure persists across collapse, it enters the next antecedent configuration and thereby biases subsequent reachable geometry toward continuations compatible with that realised structure. Momentum is consequently not introduced by an unexplained grounding arrow from Narrative to a second ontological object. It is the directional dynamical consequence of operative Narrative in a geometry that supplies direction.

The resulting law is structurally adjacent to established work on hysteresis, path dependence, belief revision, predictive coding, motor schemas, and institutional lock-in, but is not identical to any one of them. The contribution here is the common frame-relative formulation: realised history becomes first-class state, present state carries only the history that remains causally operative, and directional persistence follows from its presence in the next possibility relation.
\end{abstract}

\section{Dependencies and scope}\label{sec:dependencies}
Graphic Equalisation supplies deterministic realised transition
\[
G_t\xrightarrow{\mathcal C_t}G_{t+1},
\]
and the Seldon-manifold theory supplies counterfactual reachable geometry
\[
\Seldon_{h,\rho}(\Frame_t).
\]
This paper adds a measurable higher-frame quantity carrying realised history. It does not require cognition, self-models, language, or deliberative agency.

The nearest established structures are history-dependent dynamics and hysteresis, in which present response depends on retained information about past input \citep{mayergoyz2003}; economic and institutional path dependence, in which earlier events alter later reachable trajectories \citep{david1985,arthur1989,pierson2000}; formal belief revision, in which changing one proposition may require revision of a larger dependency structure \citep{agm1985}; and predictive-processing models in which present expectations shape interpretation of subsequent input \citep{rao1999}. These provide prior art for particular behaviours. The present claim is the cross-domain identity of the underlying frame-relative structure.

\section{Narrative as a first-class dimension}\label{sec:narrative}
Let $W$ individuate graph $G$.

\begin{definition}[Narrative dimension]\label{def:narrative-dimension}
The Narrative dimension $D_{\Narr,W:G}$ exists in frame $W$ when $W$ represents a history-bearing relation concerning $G$ whose value is causally operative in subsequent transition or reachability structure.
\end{definition}

Before that relation exists,
\[
D_{\Narr,W:G}=\operatorname{NULL},
\]
and $\Narr_{W:G,t}$ is undefined. This is not the numerical state $\Narr_{W:G,t}=0$. Existence of the dimension precedes value on the dimension.

\begin{definition}[Narrative]\label{def:narrative}
Where $D_{\Narr,W:G}$ exists,
\[
\boxed{\Narr_{W:G,t}\in\mathbb Z}
\]
is the discrete frame-relative measure of realised-history structure concerning $G$ that remains represented and operative in $W$ at realised step $t$.
\end{definition}

The integer is a state index rather than a globally fixed unit. A grounding may map it to arbitrary dimensional scales or dynamically varying quantisation. A value of zero is a legitimate value on an instantiated Narrative dimension; it is not absence of that dimension.

\begin{proposition}[No emergence gap]\label{prop:no-emergence-gap}
The absence of Narrative as a primitive in $G=(g,r,\Gamma)$ does not make Narrative non-physical or metaphorical. Where the defining relation exists in $W$, Narrative is ordinary first-class state in $W$.
\end{proposition}

This is a direct use of the higher-frame principle established in Graphic Equalisation: a relation may be absent from the lower graph while being fully real in a frame containing the relation required to define it.

\section{Narrative and realised history}\label{sec:history}
Realised history is the represented succession
\[
\Hist_{W:G,\le t}=G_0\xrightarrow{\mathcal C_0}G_1\xrightarrow{\mathcal C_1}\cdots\xrightarrow{\mathcal C_{t-1}}G_t.
\]
Narrative is not necessarily the complete history. It is the currently operative history-bearing structure:
\[
\Hist_{W:G,\le t}\longrightarrow\Narr_{W:G,t}.
\]

\begin{proposition}[Historical sufficiency of present structure]\label{prop:historical-sufficiency}
Past events affect future dynamics only insofar as their consequences remain represented in present graph or boundary structure.
\end{proposition}

This proposition is closely related to the state-variable treatment of systems with memory: the past influences future behaviour through whichever memory variables remain operative in the present representation. Hysteresis theory makes this explicit in systems whose present response retains selective information about past extrema or branches \citep{mayergoyz1991,mayergoyz2003}. The present theory generalises the statement across frames and graph types.

Narrative is therefore not ``everything that happened.'' It is the part of realised history that has not disappeared dynamically.

\section{Direction and momentum}\label{sec:momentum}
Narrative magnitude alone is not a vector. Direction exists only where the relevant frame also instantiates a directional dimension or ordering over continuation.

Let $\hat{\mathbf d}_{W:G,t}$ denote a locally represented direction associated with the history-bearing relation. In a grounding where scalar Narrative magnitude and directional coordinates are commensurable, a convenient coordinate representation is
\[
\boxed{\mathbf p_{W:G,t}=\phi_W(\Narr_{W:G,t})\,\hat{\mathbf d}_{W:G,t},}
\]
where $\phi_W$ supplies the frame-specific scale. The special case $\phi_W(n)=n$ gives $\mathbf p=\Narr\hat{\mathbf d}$.

\begin{definition}[Narrative momentum]\label{def:narrative-momentum}
Narrative momentum is the directional persistence in reachable future geometry induced by operative Narrative when a directional dimension exists in the frame.
\end{definition}

Momentum is therefore derived from first-class Narrative plus represented geometry; it is not posited as a second unexplained substance. Classical momentum is one exact grounding of this structural form, not the definition from which the abstract quantity is borrowed.

\section{The Law of Narrative Momentum}\label{sec:law}
\begin{hypothesis}[Law of Narrative Momentum]\label{hyp:narrative-momentum}
Let the Narrative dimension exist in frame $W$, let $\Narr_{W:G,t}$ be operative, and let the relations supporting that Narrative remain materially unchanged across the next realised transition. Then subsequent reachable geometry retains a directional component aligned with the continuation represented by that Narrative. Material redirection requires an interaction that changes the relations carrying the operative history.
\end{hypothesis}

The key recursion is
\[
\boxed{\text{history survives as structure}\Rightarrow\text{history enters the next possibility relation}.}
\]

The law itself is deterministic. For empirical work over repeated cases where the observer does not represent every determining dimension, an ensemble statistic such as
\[
\frac{1}{m}\sum_{i=1}^{m}\mathbf p^{(i)}_{t+1}\cdot\widehat{\mathbf p}^{(i)}_t>0
\]
may test directional persistence across cases. The statistic measures the observer's population-level evidence; it does not add stochasticity to the Graph Law.

\section{Narrative revision}\label{sec:revision}
\begin{definition}[Narrative revision]\label{def:narrative-revision}
A material Narrative revision is a realised change
\[
\Delta\Narr\neq0
\]
large enough to alter future distinctions or direction at the operative frame and resolution.
\end{definition}

A proposition may change while Narrative does not. Conversely, a small local change may induce a large Narrative revision when it alters a highly connected history-bearing relation.

Formal belief-revision theory already distinguishes local proposition change from revision of a larger belief set and studies the conditions under which contraction or revision preserves as much prior structure as possible \citep{agm1985}. The present framework differs by treating the relevant dependency structure as dynamical graph state rather than purely propositional closure, but the structural predecessor is direct.

\section{Structural lock-in}\label{sec:lockin}
\begin{hypothesis}[Structural lock-in]\label{hyp:lockin}
For fixed frame and revision class, the minimum intervention required for material Narrative revision tends to increase with the amount of present relational structure that must be modified, bypassed, or rendered irrelevant.
\end{hypothesis}

If $K_R(\Narr)$ is a discrete measure of Narrative-supporting structure, an appropriate monotone grounding predicts that increasing $K_R$ does not decrease minimum revision cost after relevant capacity and geometry are controlled.

This generalises familiar forms of lock-in. In technology and economics, increasing returns can make an historically selected path progressively harder to leave \citep{arthur1989,david1985}. In institutional analysis, timing, sequence, and increasing returns can similarly make established arrangements difficult to reverse \citep{pierson2000}. These are groundings of the structural-lock-in hypothesis, not its proof.

\section{History-sensitive projection}\label{sec:history-projection}
Let $H_t$ be dynamically relevant history and $\widehat H_t$ represented history.

\begin{hypothesis}[History-sensitive projection]\label{hyp:history-projection}
Where two represented histories imply different operative present relations despite selected instantaneous observables matching, their projected reachable geometries diverge over some relevant horizon:
\[
\widehat H_t\neq H_t
\Rightarrow
\widehat\Seldon_{h,\rho}(\Frame_t)\neq\Seldon_{h,\rho}(\Frame_t).
\]
\end{hypothesis}

The load-bearing clause is ``different operative present relations.'' If the complete causally operative present state is identical, Graphic Equalisation is deterministic and the histories cannot produce different successors merely because they were different histories. History matters only where it survives in state, exactly as Proposition~\ref{prop:historical-sufficiency} requires.

Predictive-coding models provide a cognitive analogue: higher-level predictive state shapes the processing of subsequent input, while residual error updates that state \citep{rao1999}. Motor-control schema theory likewise treats prior movement experience as modifying generalized structures used to parameterise later action \citep{schmidt1975}. Neither is claimed to be identical to Narrative Momentum; both demonstrate established domain-specific cases in which retained history changes future projection.

\section{Mechanical grounding}\label{sec:mechanics}
Classical momentum is a constrained quantitative grounding. In a mechanical frame, realised positional history together with a time embedding determines velocity; mass supplies the established mechanical scale and
\[
\mathbf p=m\mathbf v.
\]
The present theory does not redefine classical momentum. It identifies the same abstract ingredients---frame-relative history-bearing state, direction, persistence, and redirection---in a domain where the quantitative law is already exact.

The analogy must not be overread. Narrative is not asserted to have SI momentum units, and non-mechanical groundings need not obey conservation of mechanical momentum.

\section{Cognitive and institutional groundings}\label{sec:groundings}
In cognition, Narrative-supporting relations may include expectations, habits, motor programmes, autobiographical commitments, causal explanations, or identity-bearing relations. Predictive coding provides an established example in which persistent higher-level expectations shape subsequent interpretation \citep{rao1999}; schema theory supplies a motor-control example in which learned generalized structures condition future movement \citep{schmidt1975}; and formal belief revision provides a model of dependency-sensitive restructuring \citep{agm1985}.

In institutions, contracts, procedures, staffing, compatibility commitments, infrastructure, capital allocation, and reputation can carry Narrative. Path-dependence research provides direct precedent for the claim that historically accumulated structures can constrain later possibilities and become progressively costly to reverse \citep{arthur1989,pierson2000}. Routine output may therefore remain stable even when the institution's capacity to materially revise the structure carrying that routine has degraded.

These are groundings of the abstract structure, not claims that cognition or institutions are literally Newtonian particles.

\section{Relationship to Narrativium}\label{sec:narrativium}
Narrative Momentum describes persistence. Narrativium describes capacity for deformation:
\[
\boxed{\Narr=\text{operative history-bearing structure},\qquad N=\text{reachable capacity to change structure}.}
\]
Revision horizons produced by comparing current capacity with minimum revision cost are developed separately.

\section{Positioning and prior art}\label{sec:prior-art}
The core ingredients of Narrative Momentum are individually old. Hysteresis theory formalises selective memory in dynamical systems \citep{mayergoyz1991,mayergoyz2003}. Path dependence and increasing returns formalise historically conditioned lock-in in economics and institutions \citep{david1985,arthur1989,pierson2000}. AGM theory formalises dependency-sensitive belief revision \citep{agm1985}. Predictive coding and motor-schema theories formalise history-bearing internal structure affecting future processing and action \citep{rao1999,schmidt1975}.

The claimed contribution is the reduction rather than the ingredients: these phenomena can be expressed as instances of one frame-relative pattern,
\[
\boxed{\text{realised history}\rightarrow\text{operative present structure}\rightarrow\text{directional persistence of reachable futures}.}
\]
Graphic Equalisation supplies the causal reason the pattern works: only present structure enters the next deterministic transition relation. Seldon manifolds supply the geometry in which that persistence is expressed. Narrative names the measurable higher-frame state that carries the surviving history.

\section{Falsifiability commitments}\label{sec:falsifiability}
\begin{enumerate}[leftmargin=*]
\item \textbf{First-class Narrative.} Fails for a claimed grounding if Narrative cannot be measured as frame-level state and exists only as an after-the-fact interpretive label.
\item \textbf{Dimension/value typing.} Fails if absence of the Narrative relation must be represented as a numerical value of Narrative rather than as absence of the dimension.
\item \textbf{Discrete Narrative.} Fails if Narrative-dependent observable behaviour requires a non-discrete primitive that no dynamically quantised discrete representation can preserve.
\item \textbf{Historical sufficiency.} Fails if erased history with no present causal trace systematically changes future transition structure.
\item \textbf{Narrative Momentum.} Fails for a specified system class if operative Narrative-supporting structure and its represented direction persist while the associated directional tendency repeatedly vanishes or reverses without a relevant interaction.
\item \textbf{Structural lock-in.} Fails if increasing the structure that must be revised systematically lowers minimum revision cost after relevant capacity and geometry are controlled.
\item \textbf{History-sensitive projection.} Fails if histories that produce different operative present relations never produce different reachable projections in the claimed domain; conversely, the hypothesis is misapplied if different histories with identical causally complete present state are claimed to produce different deterministic successors.
\end{enumerate}

\section{Conclusion}\label{sec:conclusion}
Narrative is not a metaphor smuggled into dynamics. It is first-class state in the frame that contains the relations from which it is measured:
\[
\boxed{\text{collapse}\rightarrow\text{realised history}\rightarrow\Narr\rightarrow\text{directional persistence in reachable geometry}.}
\]
Momentum no longer requires a magical bridge. Narrative is already dynamical state in the higher frame; Narrative momentum is the directional persistence produced when that state participates in subsequent reachable geometry.

The constituent ideas have long precedents. The proposed novelty is that the same minimal recursive structure accounts for their common form without making Narrative a primitive of every graph or treating history as a force that acts from the past without surviving present structure.

\begin{thebibliography}{99}
\bibitem{mayergoyz1991}
I. D. Mayergoyz.
\newblock \emph{Mathematical Models of Hysteresis}.
\newblock Springer, New York, 1991.

\bibitem{mayergoyz2003}
I. D. Mayergoyz.
\newblock \emph{Mathematical Models of Hysteresis and Their Applications}.
\newblock Academic Press, 2nd ed., 2003.

\bibitem{david1985}
P. A. David.
\newblock Clio and the economics of QWERTY.
\newblock \emph{American Economic Review}, 75(2):332--337, 1985.

\bibitem{arthur1989}
W. B. Arthur.
\newblock Competing technologies, increasing returns, and lock-in by historical events.
\newblock \emph{The Economic Journal}, 99(394):116--131, 1989.

\bibitem{pierson2000}
P. Pierson.
\newblock Increasing returns, path dependence, and the study of politics.
\newblock \emph{American Political Science Review}, 94(2):251--267, 2000.

\bibitem{agm1985}
C. E. Alchourr\'on, P. G\"ardenfors, and D. Makinson.
\newblock On the logic of theory change: Partial meet contraction and revision functions.
\newblock \emph{Journal of Symbolic Logic}, 50(2):510--530, 1985.

\bibitem{rao1999}
R. P. N. Rao and D. H. Ballard.
\newblock Predictive coding in the visual cortex: a functional interpretation of some extra-classical receptive-field effects.
\newblock \emph{Nature Neuroscience}, 2:79--87, 1999.

\bibitem{schmidt1975}
R. A. Schmidt.
\newblock A schema theory of discrete motor skill learning.
\newblock \emph{Psychological Review}, 82(4):225--260, 1975.
\end{thebibliography}

\end{document}
